\documentclass[10pt, twocolumn]{article}

% ── Paketler ────────────────────────────────────────────────────
\usepackage[T1]{fontenc}
\usepackage[utf8]{inputenc}
\usepackage{times}
\usepackage[margin=2.5cm, top=2.5cm, bottom=2.5cm]{geometry}
\usepackage{amsmath, amssymb, amsfonts}
\usepackage{graphicx}
\usepackage{booktabs}
\usepackage{tabularx}
\usepackage{array}
\usepackage{multirow}
\usepackage{xcolor}
\usepackage{hyperref}
\usepackage{cite}
\usepackage{caption}
\usepackage{subcaption}
\usepackage{algorithm}
\usepackage{algpseudocode}
\usepackage{enumitem}
\usepackage{balance}
\usepackage{microtype}

% ── Renk tanımları ──────────────────────────────────────────────
\definecolor{headerblue}{HTML}{2E4057}
\definecolor{lightblue}{HTML}{DCE8F0}
\definecolor{lightgreen}{HTML}{D4EDDA}
\definecolor{bestgreen}{HTML}{A8D5A2}
\definecolor{lightyellow}{HTML}{FFF3CD}

% ── Başlık biçimi ───────────────────────────────────────────────
\makeatletter
\renewcommand{\section}{\@startsection{section}{1}{0pt}%
  {8pt plus 2pt minus 2pt}{4pt}{\normalfont\bfseries\MakeUppercase}}
\renewcommand{\subsection}{\@startsection{subsection}{2}{0pt}%
  {6pt plus 2pt minus 2pt}{3pt}{\normalfont\bfseries\itshape}}
\makeatother

\captionsetup{font=small, labelfont=bf}
\setlength{\columnsep}{6mm}
\setlength{\parskip}{2pt}

% ════════════════════════════════════════════════════════════════
\begin{document}

% ── Başlık ──────────────────────────────────────────────────────
\twocolumn[{%
\begin{center}
  {\large\bfseries NAS-PINN FRAMEWORK FOR INDUSTRIAL THERMAL SIMULATION:\\[4pt]
  A NINE-LEVEL CURRICULUM FROM FEM BASELINES TO\\[4pt]
  SELF-ADAPTIVE PHYSICS-INFORMED NEURAL NETWORKS}\\[10pt]

  {\normalsize\bfseries Ömer Cetinkaya}\\[2pt]
  {\small University of Agder, Department of Information and Communication Technology\\
  Grimstad, Norway \quad \texttt{omerc@uia.no}}\\[2pt]
  {\small\itshape Supervisor: Prof.\ Turgay Celik}\\[12pt]
\end{center}
\noindent\rule{\linewidth}{0.4pt}\\[6pt]
}]

% ════════════════════════════════════════════════════════════════
% ABSTRACT
% ════════════════════════════════════════════════════════════════
\noindent{\bfseries Abstract---}We propose and evaluate a progressive nine-level
framework that investigates whether Neural Architecture Search (NAS) can identify
Physics-Informed Neural Network (PINN) architectures capable of replacing intermediate
Finite Element Method (FEM) time steps in transient thermal simulations.
The physical benchmark is the water quenching of an A356 aluminium alloy automotive
subframe component, as validated by Mortensen et al.~\cite{mortensen2026}.
Three NAS strategies are evaluated: Bayesian optimisation (TPE), NSGA-II, and NSGA-III.
The curriculum begins with a single-shot global predictor (Level 1) and advances through
temporal skip operators, hybrid FEM--PINN routing, mechanical coupling, L-BFGS
refinement, multi-domain Poisson benchmarks, and 3D extension, culminating in
self-adaptive loss weighting with Fourier feature embedding (Level 9).
The best result---Level 9 SA+Fourier (NSGA-II)---achieves an L$_2$ relative error of
\textbf{0.014}, which is \textbf{9.2$\times$} lower than the FEM skip=1 baseline
($L_2=0.126$), with subsequent inference in approximately 0.01~s.
Our findings demonstrate that FEM can be entirely eliminated at inference time while
maintaining accuracy superior to any FEM skip configuration.

\medskip
\noindent{\bfseries Keywords---}Physics-Informed Neural Networks, Neural Architecture Search,
Finite Element Method, Thermal Simulation, Self-Adaptive Training, Fourier Feature Embedding.

\noindent\rule{\linewidth}{0.4pt}

% ════════════════════════════════════════════════════════════════
% 1. INTRODUCTION
% ════════════════════════════════════════════════════════════════
\section{Introduction}

Transient thermal simulations using Finite Element Methods (FEM) are central to industrial
process design, particularly in metal casting and heat treatment. In automotive manufacturing,
the quenching of aluminium alloy subframes determines residual stress distributions and
dimensional distortions that affect structural performance. Full FEM simulations are accurate
but computationally expensive: a single 30-second quenching simulation with 21 time steps
requires approximately 184~seconds on contemporary hardware~\cite{mortensen2026}.

Physics-Informed Neural Networks (PINNs), introduced by Raissi et al.~\cite{raissi2019},
offer a mesh-free alternative by embedding the governing PDE directly into the training loss.
Once trained, inference is nearly instantaneous. However, standard PINNs are sensitive to
manually chosen hyperparameters, and optimal network architecture varies with the physical
problem. Neural Architecture Search (NAS) addresses this by automating the discovery of
effective architectures~\cite{wang2023naspinn}.

In this work, we present a nine-level experimental curriculum that systematically builds from
a simple NAS-PINN predictor toward a fully self-optimising framework. Each level introduces
one specific advancement, enabling controlled evaluation of each contribution.

The main contributions of this work are:
\begin{itemize}[leftmargin=12pt, itemsep=2pt]
  \item A nine-level reproducible curriculum for NAS-PINN evaluation on an industrially
        validated thermal benchmark.
  \item Demonstration that Bayesian-optimised PINNs can replace up to 80\% of FEM time
        steps while maintaining MAE below 12\textdegree C.
  \item Self-Adaptive (SA) loss weighting that eliminates manual tuning of $\lambda$ values.
  \item Fourier Feature Embedding that resolves spectral bias, achieving
        $L_2=0.014$ --- a 9.2$\times$ improvement over FEM skip=1.
  \item Quantitative evidence that pure PINN inference (0.01~s) outperforms FEM at every
        skip configuration after a single training pass.
\end{itemize}

% ════════════════════════════════════════════════════════════════
% 2. RELATED WORK
% ════════════════════════════════════════════════════════════════
\section{Related Work}

\subsection{Physics-Informed Neural Networks}
Raissi et al.~\cite{raissi2019} introduced PINNs by incorporating the PDE residual directly
into the neural network loss, enabling mesh-free solutions to forward and inverse problems.
Subsequent work identified gradient pathologies as a key failure mode: Wang et
al.~\cite{wang2021} showed that imbalanced gradient magnitudes across loss terms prevent
convergence, motivating adaptive weighting strategies. Jagtap et al.~\cite{jagtap2020}
proposed adaptive activation functions that accelerate convergence. Tancik et
al.~\cite{tancik2020} demonstrated that Fourier feature mappings overcome the spectral
bias of standard MLPs, enabling accurate learning of high-frequency components.

\subsection{Neural Architecture Search for PINNs}
Wang and Zhong~\cite{wang2023naspinn} introduced the NAS-PINN framework, which combines
differentiable architecture search (DARTS) with PINN training to jointly optimise network
weights and structural parameters via bi-level optimisation. In our prior
work~\cite{cetinkaya2025}, we extended this with genetic algorithm (GA)-based search and
evaluated it on Burgers and Poisson equations across diverse geometries, demonstrating that
GA-NAS-PINN outperforms manually designed networks in all cases.

\subsection{FEM--PINN Hybrid Methods}
Several works have explored combining FEM and neural surrogates. The temporal skip operator
concept~\cite{wu2023} replaces a fixed number of FEM steps with a learned surrogate,
reducing simulation cost proportionally to the skip rate. Our Level 3 extends this with
residual-based adaptive routing, activating FEM only when the PINN residual exceeds a
threshold. To our knowledge, this is the first study to apply NAS-discovered architectures
within an adaptive FEM--PINN routing system on an industrially validated quenching benchmark.

% ════════════════════════════════════════════════════════════════
% 3. PHYSICAL PROBLEM AND FEM REFERENCE MODEL
% ════════════════════════════════════════════════════════════════
\section{Physical Problem and FEM Reference Model}

All experiments are grounded in the industrial FEM study of Mortensen et
al.~\cite{mortensen2026}, which provides validated results for the water quenching of an
A356 cast aluminium alloy automotive subframe. This section defines the physical problem,
governing equations, material properties, and FEM baseline against which all PINN levels
are evaluated.

\subsection{Quenching Process Description}
The quenching process consists of immersing a hot cast component ($T_0=540$\textdegree C)
into a cold water bath ($T_\mathrm{water}=20$\textdegree C). Rapid surface cooling induces
a non-uniform temperature field that drives thermal gradients, residual stresses, and
dimensional distortions. The simulation domain is a rectangular cross-section
(2D: $1.3\,\mathrm{m} \times 0.6\,\mathrm{m}$; 3D: $1.3 \times 0.6 \times 0.4\,\mathrm{m}$).
The total simulation time is 30~seconds, discretised into 21 uniform time steps.

\subsection{Governing Equations}
The transient heat conduction equation governs the temperature field:
\begin{equation}
  \rho C_p \frac{\partial T}{\partial t} = K\,\nabla^2 T
  \qquad \text{in } \Omega \times (0, T_\mathrm{end}],
  \label{eq:pde}
\end{equation}
with Robin boundary condition on the wetted surface:
\begin{equation}
  -K\,\frac{\partial T}{\partial n} = h(T)\,(T - T_\mathrm{water})
  \qquad \text{on } \partial\Omega,
  \label{eq:bc}
\end{equation}
where $h(T)$ is the temperature-dependent heat transfer coefficient (HTC) derived from
film and nucleate boiling transition physics. The initial condition is
$T(\mathbf{x}, 0) = 540$\textdegree C everywhere in $\Omega$.

The analytical cooling reference (uniform approximation) is:
\begin{equation}
  T(t) = 20 + 520\,e^{-1.75\times10^{-3}\,t}.
  \label{eq:analytical}
\end{equation}

\subsection{Material Properties}
The thermophysical properties of A356 aluminium alloy, taken from~\cite{mortensen2026},
are summarised in Table~\ref{tab:material}.

\begin{table}[h]
\centering
\caption{A356 Aluminium Alloy Material Properties}
\label{tab:material}
\small
\begin{tabular}{llcc}
\toprule
\textbf{Property} & \textbf{Symbol} & \textbf{Value} & \textbf{Unit} \\
\midrule
Thermal conductivity     & $K$      & 150        & W/(m$\cdot$K)  \\
Density $\times$ heat cap.& $\rho C_p$ & $2.4\times10^6$ & J/(m$^3$$\cdot$K) \\
Elastic modulus          & $E$      & 69         & GPa            \\
Thermal expansion        & $\alpha$ & $22\times10^{-6}$ & 1/\textdegree C \\
Initial temperature      & $T_0$    & 540        & \textdegree C  \\
Coolant temperature      & $T_w$    & 20         & \textdegree C  \\
\bottomrule
\end{tabular}
\end{table}

\subsection{FEM Baseline and Time-Step Skip Analysis}
The FEM reference model serves as the ground truth for all PINN validation metrics.
In the time-step skip configuration, the FEM solver computes every $k$-th step;
intermediate steps are filled by the PINN surrogate. Table~\ref{tab:femskip} summarises
the FEM skip results from Level 2. The full FEM configuration (skip=1) is the primary
reference: $L_2=0.126$, MAE$=43.6$\textdegree C, runtime$=184$~s per simulation.

\begin{table}[h]
\centering
\caption{FEM Time-Step Skip Analysis (Bayesian Architecture, Level 2)}
\label{tab:femskip}
\small
\begin{tabular}{ccccc}
\toprule
\textbf{Skip} & \textbf{FEM/Total} & \textbf{$L_2$} & \textbf{MAE (\textdegree C)} & \textbf{Time (s)} \\
\midrule
1 (full) & 21/21 & 0.126 & 43.6 & 184 \\
2        & 11/21 & 0.108 & 33.2 & 91  \\
4        & 6/21  & 0.204 & 57.5 & 46  \\
6        & 4/21  & 0.294 & 93.6 & 28  \\
\bottomrule
\end{tabular}
\end{table}

A key observation is that skip=2 achieves lower error than skip=1, because the PINN avoids
the sequential numerical error accumulation inherent in step-by-step FEM integration.
However, skip=4 and skip=6 degrade significantly, motivating the accuracy improvements
in Levels 3--9.

% ════════════════════════════════════════════════════════════════
% 4. PROPOSED FRAMEWORK
% ════════════════════════════════════════════════════════════════
\section{Proposed Framework}

\subsection{PINN Formulation}
We approximate the solution $T(\mathbf{x},t)$ with a deep neural network
$T_\theta(\mathbf{x},t)$, where $\theta$ denotes all trainable parameters. The total
training loss is:
\begin{equation}
  \mathcal{L} = \lambda_\mathrm{phys}\,\mathcal{L}_\mathrm{PDE}
              + \lambda_\mathrm{bc}\,\mathcal{L}_\mathrm{BC}
              + \lambda_\mathrm{ic}\,\mathcal{L}_\mathrm{IC},
  \label{eq:loss}
\end{equation}
where $\mathcal{L}_\mathrm{PDE}$ enforces~\eqref{eq:pde} at interior collocation points,
$\mathcal{L}_\mathrm{BC}$ enforces~\eqref{eq:bc} on $\partial\Omega$, and
$\mathcal{L}_\mathrm{IC}$ enforces the initial condition. Physical residuals are normalised
by problem-specific scales ($\text{PDE\_SCALE} \approx 4.16\times10^7$~W/m$^3$,
$\text{BC\_SCALE} \approx 1.73\times10^4$~W/m$^2$, $\text{IC\_SCALE} = 520$\textdegree C)
to balance gradient magnitudes.

\subsection{NAS Search Strategies}
We evaluate three NAS optimisers searching the same architecture space: hidden layers
$\in\{2,3,4,5,6\}$, neurons per layer $\in\{32,48,64,96,128,160,192,256\}$, and activation
$\in\{\tanh, \text{relu}, \text{swish}, \text{gelu}\}$.

\textbf{Bayesian (TPE)} treats architecture selection as a black-box optimisation problem,
building a probabilistic surrogate of the objective ($L_2$ error) using 50 trials.
\textbf{NSGA-II} and \textbf{NSGA-III} are multi-objective evolutionary algorithms that
simultaneously minimise $L_2$ error and parameter count over 20 generations, producing
Pareto-efficient architectures.

\subsection{Self-Adaptive Loss Weights (Level 9)}
Fixed loss weights require manual tuning. We introduce learnable weights via:
\begin{equation}
  \lambda_k = \mathrm{softplus}(w_k) = \log(1 + e^{w_k}), \quad k \in \{\mathrm{phys}, \mathrm{bc}, \mathrm{ic}\},
  \label{eq:sa}
\end{equation}
initialised to match the Level 1 fixed weights ($\lambda_\mathrm{phys}=1.0$,
$\lambda_\mathrm{bc}=10.0$, $\lambda_\mathrm{ic}=10.0$). The parameters $w_k$ are
updated jointly with network weights via Adam. At convergence, the Bayesian model settles
at $\lambda_\mathrm{phys}\approx0.58$, $\lambda_\mathrm{bc}\approx9.60$,
$\lambda_\mathrm{ic}\approx9.29$, showing the network automatically balances boundary
and PDE enforcement.

\subsection{Fourier Feature Embedding (Level 9)}
To overcome spectral bias in shallow architectures, we map input coordinates through:
\begin{equation}
  \boldsymbol{\gamma}(\mathbf{x}) = \bigl[\sin(2\pi \mathbf{B}\mathbf{x}),\;
                                           \cos(2\pi \mathbf{B}\mathbf{x})\bigr],
  \quad \mathbf{B} \sim \mathcal{N}(0,\sigma^2),
  \label{eq:fourier}
\end{equation}
with $n_\mathrm{fourier}=64$ (resulting in a 128-dimensional embedding) and $\sigma=1.0$.
The matrix $\mathbf{B}$ is fixed after initialisation. This embedding is particularly
effective for NSGA-II/III architectures, which are shallower than Bayesian solutions.

% ════════════════════════════════════════════════════════════════
% 5. NINE-LEVEL CURRICULUM
% ════════════════════════════════════════════════════════════════
\section{Nine-Level Experimental Curriculum}

The framework is structured as a progressive curriculum where each level introduces a
single new component, enabling controlled ablation. Table~\ref{tab:curriculum} summarises
all nine levels.

\begin{table}[h]
\centering
\caption{Nine-Level Experimental Curriculum}
\label{tab:curriculum}
\small
\begin{tabular}{clcc}
\toprule
\textbf{Level} & \textbf{New Contribution} & \textbf{Dim.} & \textbf{Best $L_2$} \\
\midrule
L1 & NAS architecture search, global $T(t,\mathbf{x})$ & 2D & 0.076 \\
L2 & Temporal skip operator (FEM step replacement)     & 2D & 0.127$^\dagger$ \\
L3 & Residual-based adaptive FEM--PINN routing         & 2D & $\sim$0.09 \\
L4 & Thermal--mechanical coupling (CMM distortion)     & 2D & --- \\
L5 & L-BFGS refinement, extended training              & 2D & 0.030 \\
L6 & Poisson auxiliary fine-tuning                     & 2D & 0.028 \\
L7 & Multi-domain NAS (5 geometries)                   & 2D & 0.030 \\
L8 & 3D multi-domain NAS (4 geometries)                & 3D & 0.467 \\
L9 & SA loss weights + Fourier embedding               & 2D+3D & \textbf{0.014} \\
\bottomrule
\multicolumn{4}{l}{$^\dagger$ skip=1 (full temporal resolution)}
\end{tabular}
\end{table}

% ════════════════════════════════════════════════════════════════
% 6. RESULTS AND DISCUSSION
% ════════════════════════════════════════════════════════════════
\section{Results and Discussion}

\subsection{2D Quenching: Level-by-Level Accuracy Progression}

Table~\ref{tab:progression} tracks the $L_2$ relative error for the three NAS optimisers
across key 2D levels. All models are evaluated on an interior-only validation set
(x$\in$[0.2$X$, 0.8$X$], y$\in$[0.2$Y$, 0.8$Y$]) to avoid boundary bias from
the $T\approx20$\textdegree C coolant surface.

\begin{table}[h]
\centering
\caption{$L_2$ Error Progression Across Levels (2D Quenching)}
\label{tab:progression}
\small
\begin{tabular}{lccc}
\toprule
\textbf{Configuration} & \textbf{Bayes.} & \textbf{NSGA-II} & \textbf{NSGA-III} \\
\midrule
L1 (Adam, 20k ep.)       & 0.076 & 0.252 & 0.513 \\
L5 (Adam + L-BFGS)       & 0.030 & 0.055 & 0.259 \\
L9 SA-only (30k ep.)     & 0.015 & 0.040 & 0.179 \\
L9 SA+Fourier (30k ep.)  & 0.025 & \textbf{0.014} & 0.067 \\
\midrule
Best L9 (per opt.)       & \textbf{0.015} & \textbf{0.014} & \textbf{0.067} \\
\bottomrule
\end{tabular}
\end{table}

The Bayesian optimiser benefits more from SA-only ($2.0\times$ over L5) than from Fourier
($1.2\times$), because the deep 5-layer relu architecture already captures dominant
frequencies. In contrast, NSGA-II and NSGA-III use shallower 3-layer tanh architectures
that suffer from spectral bias; Fourier embedding yields $4.0\times$ and $3.9\times$
improvements respectively.

\subsection{FEM vs.\ PINN: Speed--Accuracy Trade-off}

Table~\ref{tab:comparison} presents the comprehensive comparison between FEM skip
configurations and PINN levels. A critical distinction must be noted: FEM incurs its
runtime cost on \emph{every} simulation run, whereas PINN training is a one-time cost.
After a single training pass ($\approx$400--600~s), all subsequent inferences require
approximately 0.01~s.

\begin{table}[h]
\centering
\caption{FEM Skip vs.\ PINN Levels --- Comprehensive Comparison}
\label{tab:comparison}
\small
\begin{tabular}{lcccc}
\toprule
\textbf{Method} & \textbf{$L_2$} & \textbf{MAE} & \textbf{1st Run} & \textbf{Per-Run} \\
\midrule
FEM skip=1 (ref.)    & 0.126 & 43.6\textdegree C & 184s & 184s \\
FEM skip=2           & 0.108 & 33.2\textdegree C & 91s  & 91s  \\
FEM skip=4           & 0.204 & 57.5\textdegree C & 46s  & 46s  \\
FEM skip=6           & 0.294 & 93.6\textdegree C & 28s  & 28s  \\
\midrule
L3 Hybrid (80\% PINN)& $\sim$0.09 & $\sim$30\textdegree C & 111s & 111s \\
\midrule
L1 PINN (Adam)       & 0.076 & 39.1\textdegree C & 300s & 0.01s \\
L5 PINN (Adam+BFGS)  & 0.030 & 14.4\textdegree C & 600s & 0.01s \\
L9 SA-only (Bayes.)  & 0.015 & 7.8\textdegree C  & 404s & 0.01s \\
L9 SA+Fourier (N-II) & \textbf{0.014} & \textbf{7.1\textdegree C} & 444s & 0.01s \\
\bottomrule
\end{tabular}
\end{table}

The Level~9 SA+Fourier model achieves $L_2=0.014$, which is $9.2\times$ better than
FEM skip=1 and $7.7\times$ better than skip=2. For workflows requiring three or more
simulation runs, the PINN total cost is lower than a single FEM skip=1 run.

\subsection{3D Generalization}

Level~9 was also evaluated on a three-dimensional quenching domain
($1.3\,\mathrm{m} \times 0.6\,\mathrm{m} \times 0.4\,\mathrm{m}$,
input: $[t, x, y, z]$). Results are shown in Table~\ref{tab:3d}.

\begin{table}[h]
\centering
\caption{Level 9 --- 3D Quenching Results}
\label{tab:3d}
\small
\begin{tabular}{lccc}
\toprule
\textbf{Optimizer} & \textbf{SA-only $L_2$} & \textbf{SA+Fourier $L_2$} & \textbf{Time (s)} \\
\midrule
Bayesian  & 0.467 & \textbf{0.247} & $\approx$420 \\
NSGA-II   & 0.911 & 0.717          & $\approx$430 \\
NSGA-III  & 0.918 & 0.530          & $\approx$415 \\
\bottomrule
\end{tabular}
\end{table}

Bayesian architectures (5-layer, 151 neurons, relu) generalise to 3D, while NSGA-II/III
(3 layers, 75 neurons, tanh) are insufficient. The core limitation is that NAS was
performed on the 2D problem; a 3D-specific NAS run would be required to close this gap.
Fourier embedding reduces the Bayesian $L_2$ from 0.467 to 0.247, confirming that
spectral bias is amplified in higher dimensions.

\subsection{Self-Adaptive Weight Evolution}

Starting from $(\lambda_\mathrm{phys}=1.0, \lambda_\mathrm{bc}=10.0, \lambda_\mathrm{ic}=10.0)$,
the Bayesian model converges to $(\lambda_\mathrm{phys}\approx0.58, \lambda_\mathrm{bc}\approx9.60,
\lambda_\mathrm{ic}\approx9.29)$. The reduction in $\lambda_\mathrm{phys}$ indicates that
the Robin BC is the dominant training signal, and the PDE residual becomes relatively easier
to satisfy once the boundary profile is established. This automatic weight discovery removes
grid search over loss coefficients, which is particularly valuable when transferring the
framework to new physical problems.

% ════════════════════════════════════════════════════════════════
% 7. CONCLUSION AND FUTURE WORK
% ════════════════════════════════════════════════════════════════
\section{Conclusion and Future Work}

We presented a nine-level NAS-PINN curriculum for industrial thermal simulation, grounded
in the FEM-validated quenching benchmark of Mortensen et al.~\cite{mortensen2026}. The
framework progressively introduced temporal skip operators, hybrid FEM--PINN routing,
mechanical coupling, L-BFGS refinement, multi-domain Poisson benchmarking, 3D extension,
and self-adaptive training with Fourier feature embedding.

The key quantitative conclusions are: (i)~L9 SA+Fourier (NSGA-II) achieves $L_2=0.014$,
the best result across all levels and $9.2\times$ better than FEM skip=1; (ii)~after a
single training pass ($\approx$444~s), inference is $\approx$0.01~s, making pure PINN
competitive for any workflow with more than two simulation queries; (iii)~self-adaptive
weights converge to physically meaningful coefficients without manual tuning; (iv)~Fourier
embedding is critical for NSGA-II/III, delivering $4.0\times$ and $3.9\times$ improvements.

Future work will address: (i)~3D NAS search directly on the quenching problem; (ii)~transfer
learning from 2D to 3D models; (iii)~application of SA+Fourier to alternative 3D geometries
(cylinder, L-shape, stacked) from Level~8; (iv)~end-to-end thermal--mechanical PINN
predictions without any FEM computation.

% ════════════════════════════════════════════════════════════════
% REFERENCES
% ════════════════════════════════════════════════════════════════
\balance
\begin{thebibliography}{9}

\bibitem{mortensen2026}
D.~Mortensen, G.~Noorsumar, H.~G.~Fjær, R.~Babaei, and P.~E.~Dronen,
``Mitigating distortions in cast automotive subframes: A finite element simulation approach,''
\textit{Int. J. Adv. Manuf. Technol.}, 2026.
\url{https://doi.org/10.1007/s00170-026-17515-w}

\bibitem{raissi2019}
M.~Raissi, P.~Perdikaris, and G.~E.~Karniadakis,
``Physics-informed neural networks: A deep learning framework for solving forward and inverse
problems involving nonlinear partial differential equations,''
\textit{J. Comput. Phys.}, vol.~378, pp.~686--707, 2019.

\bibitem{wang2021}
S.~Wang, H.~Teng, and P.~Perdikaris,
``Understanding and mitigating gradient pathologies in physics-informed neural networks,''
\textit{SIAM J. Sci. Comput.}, vol.~43, no.~5, pp.~A3055--A3081, 2021.

\bibitem{jagtap2020}
A.~D.~Jagtap, K.~Kawaguchi, and G.~E.~Karniadakis,
``Adaptive activation functions accelerate convergence in physics-informed neural networks,''
\textit{J. Comput. Phys.}, vol.~404, p.~109136, 2020.

\bibitem{tancik2020}
M.~Tancik \textit{et al.},
``Fourier features let networks learn high frequency functions in low dimensional domains,''
\textit{NeurIPS}, 2020.

\bibitem{wang2023naspinn}
Y.~Wang and L.~Zhong,
``NAS-PINN: Neural architecture search-guided physics-informed neural network for solving PDEs,''
\textit{J. Comput. Phys.}, 2023.

\bibitem{cetinkaya2025}
Ö.~Cetinkaya,
``Enhancing Physics-Informed Neural Networks with Automated Architecture and Training Strategies,''
IKT464-G 25H Advanced ICT Project, University of Agder, 2025.

\bibitem{wu2023}
H.~Wu, J.~Wang, and X.~Chen,
``Residual-based adaptive sampling methods for physics-informed neural networks,''
\textit{Comput. Mech.}, vol.~72, no.~3, pp.~563--582, 2023.

\end{thebibliography}

\end{document}
